\begin{lemma}[Bound for fit parameters on a box]
Let $F$ be a finite-dimensional real inner-product space and let
$1\le C_0<\infty$. Write $L(K)=1+\log K$ for $K\ge2$ and
$\mathcal B_K=\{(x,t):|x|\le K,\ |t|\le K\}$. There is $C>0$, depending only on
$C_0$ and $\dim F$, such that for every $v\in\Aclass_{m,C_0}$ in the sense of
\noderef{IsFitClass} which is centered at $(0,0)$ or at $(0,1)$ with constant
$C_0$ in the sense of \noderef{IsCenteredAt}, every $K\ge2$, every
$z\in\mathcal B_K$, every $r>0$ and every fit $(p,a)$ for $v$ at $z$ of radius
$r$ with constant $C_0$ in the sense of \noderef{IsFit},
\begin{equation}\label{eq:fit-parameter-bound}
  |p|+|a|\le C\bigl(L(K)+|\log r|\bigr).
\end{equation}
\end{lemma}
\begin{proof}
Write $L(K)=1+\log K\ge1$ for $K\ge2$, and write $z_c=(0,\theta)$ with
$\theta\in\{0,1\}$ for the centering point, so that by \noderef{IsCenteredAt}
$v(z_c)=0$ and $(0,0)$ is a fit for $v$ at $z_c$ of radius $1$ with constant
$C_0$. By \noderef{IsFitClass}, $v$ admits a fit at every center of every
positive radius with constant $C_0$; in particular the hypothesis of
\noderef{FitComparisonSameCenter} holds at every center. Let $C_1>0$ be the
constant of \noderef{FitComparisonSameCenter} for this $C_0$. All constants
below depend only on $C_0$ and $\dim F$.

\emph{Step 1: a fit of a large radius at the centering point.}
Let $A\ge8$ be a constant, depending only on $C_0$ and $\dim F$, to be fixed in
Step 2, and set
\[
  R=A\,K\,L(K)\ \ \ge8\cdot2\cdot1=16>1 .
\]
By \noderef{IsFitClass} choose a fit $(p_R,a_R)$ for $v$ at $z_c$ of radius $R$.
Applying \noderef{FitComparisonSameCenter} at the center $z_c$ with the radii
$R$ and $1$ and the fits $(p_R,a_R)$ and $(0,0)$,
\[
  |p_R|+|a_R|\le C_1\bigl(1+|\log(1/R)|\bigr)=C_1\bigl(1+\log R\bigr),
\]
using $R>1$. Now $\log R=\log A+\log K+\log L(K)\le\log A+2L(K)$, since
$\log K\le1+\log K=L(K)$ and $\log L(K)\le L(K)$ (because $\log\sigma\le\sigma$
for $\sigma>0$). Hence, using $L(K)\ge1$,
\begin{equation}\label{eq:fpb-pR}
  |p_R|+|a_R|\le C_4\,L(K),
  \qquad C_4:=C_1\bigl(1+\log A+2\bigr)+1\ \ (\ge1).
\end{equation}

\emph{Step 2: the box of $(p_R,a_R)$ contains $\mathcal B_K$ with margin.}
Let $z=(x,t)$ with $|x|\le K$ and $|t|\le K$. The increment from $z_c$ to $z$ is
$(x,\,t-\theta)$, and $|t-\theta|\le K+1\le2K$. We claim
\begin{equation}\label{eq:fpb-margin}
  |t-\theta|\le\frac R4,
  \qquad \bigl|x-p_R(t-\theta)\bigr|\le\frac R4 .
\end{equation}
For the first, $|t-\theta|\le2K$ while $R/4=AKL(K)/4\ge AK/4\ge2K$ because
$A\ge8$. For the second, by \eqref{eq:fpb-pR} and $C_4L(K)\ge1$,
\[
  \bigl|x-p_R(t-\theta)\bigr|\le K+2K\,|p_R|\le K+2K\,C_4L(K)\le3C_4\,K\,L(K),
\]
and this is at most $R/4=AKL(K)/4$ as soon as $A\ge12C_4$. Since
$C_4=C_1(3+\log A)+1$ depends on $A$ only through $\log A$, the requirement
\[
  A\ge\max\bigl\{8,\ 12\bigl(C_1(3+\log A)+1\bigr)\bigr\}
\]
has the form $A\ge\alpha+\beta\log A$ with $\alpha,\beta>0$ depending only on
$C_0$ and $\dim F$, and is therefore satisfiable because $\log A/A\to0$. Fix
such an $A$ once and for all. Then \eqref{eq:fpb-pR} and
\eqref{eq:fpb-margin} hold for every $K\ge2$ and every $z\in\mathcal B_K$, with
$C_4\ge1$ depending only on $C_0$ and $\dim F$.

\emph{Step 3: fits of radius $R$ at a center in $\mathcal B_K$.}
Let $z=(x,t)\in\mathcal B_K$ and let $(p_z,a_z)$ be any fit for $v$ at $z$ of
radius $R$ with constant $C_0$. Write
$\ell_{p,a}(h,\ell)=p\cdot h+(a-|p|^2/2)\ell$, so that \noderef{IsFit} at a
center $\zeta$ and radius $\rho$ reads
$\bigl|v(\zeta+(h,\ell))-v(\zeta)-\ell_{p,a}(h,\ell)\bigr|\le C_0\rho$ on the
corresponding box.

\emph{(a) The slopes.} Let $h\in F$ with $|h|\le R/2$. The increment $(h,0)$
from $z$ lies in the box of $(p_z,a_z)$ of radius $R$, since $|0|\le R$ and
$|h-p_z\cdot0|=|h|\le R$; hence
\begin{equation}\label{eq:fpb-at-z}
  \bigl|v(x+h,t)-v(x,t)-p_z\cdot h\bigr|\le C_0R .
\end{equation}
The increments $(x,t-\theta)$ and $(x+h,t-\theta)$ from $z_c$ both lie in the box
of $(p_R,a_R)$ of radius $R$: by \eqref{eq:fpb-margin} we have
$|t-\theta|\le R/4\le R$ and $|x-p_R(t-\theta)|\le R/4\le R$, and
\[
  \bigl|x+h-p_R(t-\theta)\bigr|\le\frac R4+\frac R2\le R .
\]
Subtracting the two corresponding estimates of \noderef{IsFit}, and noting that
$\ell_{p_R,a_R}$ is linear so that its values at the two increments differ by
$\ell_{p_R,a_R}(h,0)=p_R\cdot h$, we obtain
\[
  \bigl|v(x+h,t)-v(x,t)-p_R\cdot h\bigr|\le2C_0R\qquad(|h|\le R/2).
\]
Subtracting this from \eqref{eq:fpb-at-z} gives
$|(p_z-p_R)\cdot h|\le3C_0R$ for all $|h|\le R/2$; taking
$h=(R/2)(p_z-p_R)/|p_z-p_R|$ when $p_z\ne p_R$, and dividing by $R/2$,
\begin{equation}\label{eq:fpb-slopes}
  |p_z-p_R|\le6C_0 .
\end{equation}

\emph{(b) The time coefficients.} Put $c'=\bigl(4(1+6C_0)\bigr)^{-1}$, so
$0<c'\le1/4$ and $6C_0c'\le1/4$, and consider the increment
$(p_zc'R,\,c'R)$ from $z$. It lies in the box of $(p_z,a_z)$ of radius $R$:
$|c'R|\le R$ and $|p_zc'R-p_zc'R|=0\le R$. The same space-time point, written as
an increment from $z_c$, is $\bigl(x+p_zc'R,\ t-\theta+c'R\bigr)$, and it lies
in the box of $(p_R,a_R)$ of radius $R$: by \eqref{eq:fpb-margin},
$|t-\theta+c'R|\le R/4+R/4\le R$, and by \eqref{eq:fpb-margin} and
\eqref{eq:fpb-slopes}
\[
  \bigl|x+p_zc'R-p_R(t-\theta+c'R)\bigr|
   \le\bigl|x-p_R(t-\theta)\bigr|+c'R\,|p_z-p_R|
   \le\frac R4+\frac R4\le R .
\]
Finally $(x,t-\theta)$ lies in that box by \eqref{eq:fpb-margin}.

Apply \noderef{IsFit} for $(p_z,a_z)$ at the increment $(p_zc'R,c'R)$ from $z$,
and \noderef{IsFit} for $(p_R,a_R)$ at the two increments
$(x+p_zc'R,\,t-\theta+c'R)$ and $(x,\,t-\theta)$ from $z_c$. The three
estimates involve $v$ at the two space-time points
$z+(p_zc'R,c'R)$ and $z$ only, so the values of $v$ cancel on forming the
combination (first estimate) $-$ (second) $+$ (third), and the total error is at
most $3C_0R$. Since $\ell_{p_R,a_R}$ is linear in the increment, the difference
of its values at the second and third increments is
$\ell_{p_R,a_R}(p_zc'R,c'R)$. Hence
\[
  \bigl|(\ell_{p_z,a_z}-\ell_{p_R,a_R})(p_zc'R,\,c'R)\bigr|\le3C_0R .
\]
Expanding, and using
$(p_z-p_R)\cdot p_z-\tfrac12\bigl(|p_z|^2-|p_R|^2\bigr)=\tfrac12|p_z-p_R|^2$,
\[
  (\ell_{p_z,a_z}-\ell_{p_R,a_R})(p_zc'R,c'R)
   =c'R\Bigl[a_z-a_R+\tfrac12|p_z-p_R|^2\Bigr],
\]
so dividing by $c'R>0$ and using \eqref{eq:fpb-slopes},
\begin{equation}\label{eq:fpb-times}
  |a_z-a_R|\le\frac{3C_0}{c'}+18C_0^2 .
\end{equation}

Combining \eqref{eq:fpb-pR}, \eqref{eq:fpb-slopes} and \eqref{eq:fpb-times},
and using $L(K)\ge1$, there is $C_5>0$, depending only on $C_0$ and $\dim F$,
with
\begin{equation}\label{eq:fpb-radiusR}
  |p_z|+|a_z|\le C_5\,L(K)
\end{equation}
for every $z\in\mathcal B_K$ and every fit at $z$ of radius exactly $R$.

\emph{Step 4: arbitrary radii at a center in $\mathcal B_K$.}
Let $z\in\mathcal B_K$, let $r>0$ and let $(p,a)$ be a fit for $v$ at $z$ of
radius $r$ with constant $C_0$. By \noderef{IsFitClass} choose a fit
$(p_z,a_z)$ for $v$ at $z$ of radius $R$. Applying
\noderef{FitComparisonSameCenter} at the center $z$ with the two radii $r$ and
$R$ and these two fits,
\[
  |p-p_z|+|a-a_z|\le C_1\Bigl(1+\bigl|\log(R/r)\bigr|\Bigr).
\]
Hence, with \eqref{eq:fpb-radiusR},
\[
  |p|+|a|\le C_5L(K)+C_1\Bigl(1+\bigl|\log(R/r)\bigr|\Bigr).
\]
Finally $\bigl|\log(R/r)\bigr|\le\log R+|\log r|\le\log A+2L(K)+|\log r|$ by
Step 1 and $R>1$, so with $L(K)\ge1$ there is $C>0$, depending only on $C_0$ and
$\dim F$, with
\[
  |p|+|a|\le C\bigl(L(K)+|\log r|\bigr),
\]
which is \eqref{eq:fit-parameter-bound}.
\end{proof}
